% ── SECTION 8: CONCLUSION (~600 words) ───────────────────────────────────────

\section{Conclusion: What Mutual Validation Means}
\label{sec:conclusion}

% Key points:
%   - What it means that three traditions arrived independently at the same structure.
%   - Not proof of any single tradition's authority.
%   - Evidence that all three are detecting a real feature of reality.
%   - Overlapping consensus: each tradition validates the others without dissolving.
%   - Implications for philosophy of science, science-religion dialogue, consciousness.

Three traditions built by three different civilizations, using three different
methodologies --- physical experimentation, philosophical analysis, and
spiritual discernment --- have independently arrived at the same description
of reality: that existence is constituted by relationship, that properties
emerge through relational interaction, and that the force sustaining those
relationships is, in each tradition's vocabulary, love.

The process ontology bridge transforms this observation into an argument.
Quantum mechanics does not merely resemble process ontology; it requires it,
because Bell's theorem~\citeyearpar{Bell1964} has formally eliminated every alternative.
The Bahá'í writings do not merely anticipate process ontology; they instantiate it,
forty years before \citet{Whitehead1929} named it and a century before physics
confirmed it.
Process philosophy does not merely mediate between physics and spirituality;
it names the structure that both independently discovered.

The overlapping consensus architecture preserves the integrity of each tradition.
No tradition is asked to dissolve into the others or to accept the others'
methodology as authoritative.
Each tradition validates its own thesis on its own grounds.
The convergence is the finding --- and it is more powerful than any single tradition
could produce alone.

What the convergence means for the future of science-religion dialogue, for the
philosophy of science, and for our understanding of what reality is made of, is
a question this paper opens rather than closes.
That is, perhaps, the most it can do.
And it may be enough.
